% ============================================
% SECTION 5: CIVILIZATIONAL THERMODYNAMICS
% The Hidden Energy Budget Framework
% VERSION 3.0 - FINAL HARDENED
% ============================================
% 
% FIXES APPLIED (v3.0):
% 1. Category hygiene: 3-tier taxonomy, "KIS employment" not "knowledge workers"
% 2. Dimensional hygiene: "Synthesis Investment Index" (heuristic, not metabolic)
% 3. 95% foundation: softened to "provides capacity for" not "performs"
% 4. Source hygiene: added academic anchors (Mark et al., Gonzalez & Mark)
% 5. Figures: reframed as conceptual schematics
% 6. Headlines: "pre-industrial" not "medieval produced Newton"
% ============================================

\section{Civilizational Thermodynamics: The Hidden Energy Budget}
\label{sec:civilizational}

The thermodynamic framework presented in Section~\ref{sec:methodology} explains \textit{how} to measure cognitive sovereignty. This section addresses a deeper question: \textit{why} did historical systems succeed in building sphere-capable cognition when modern systems, commanding vastly greater energy resources, systematically fail?

The answer lies in what we term the \textbf{hidden energy budget}---the constellation of conditions that enabled sustained cognitive investment, conditions so fundamental they remained invisible until their elimination revealed their necessity. The puzzle is not merely historical. It is mathematical: a civilization with 10--50 times the per-capita energy of pre-industrial Europe produces a fraction of the cognitive synthesis. Understanding this paradox requires examining not total energy but its \textit{allocation}---and specifically, the ratio of civilizational resources to active synthesizers.

\subsection{Civilizational Energy Transitions: The Smil Framework}

Before examining cognitive allocation, we must establish the material baseline. \citet{smil2017} traces civilizational energy transitions with quantitative precision, providing the empirical foundation for understanding what resources were---and are---available for cognitive investment.

Pre-industrial Western Europe operated on approximately 5--10 GJ per capita annually, barely above subsistence thresholds. The richest pre-industrial regions, benefiting from early coal and peat use alongside agricultural surplus, commanded perhaps 15--25 GJ per capita. The seventeenth-century Netherlands, often cited as the most energy-intensive pre-industrial society, consumed approximately 17 GJ per capita in peat alone, with total consumption approaching 20--25 GJ when accounting for firewood and animal labor \citep{dezeeuw1978}. By the late nineteenth century, industrial England reached approximately 100 GJ per capita. Modern affluent societies command 170--330 GJ per capita, with the European Union averaging roughly 170 GJ and the United States at the upper extreme \citep{smil2017}.

This represents an 8--17$\times$ increase in per-capita energy availability over eight centuries, depending on baseline and endpoint selection. The question this section addresses is not whether this energy increase occurred---that is documented fact---but \textit{where} that energy flows, and why so little reaches cognitive synthesis.

As Smil himself observes: ``No energetic considerations can explain the presence of Gluck, Haydn, and Mozart in the same room in Joseph II's Vienna of the 1780s'' \citep{smil2017}. Energy does not linearly produce cultural or cognitive output. The relationship is mediated by allocation structures---and it is precisely these structures that have inverted.

\subsection{The 95\% Foundation: Cognition's Hidden Architecture}

Contemporary neuroscience reveals a profound asymmetry in cognitive processing. \citet{jamadar2025} demonstrates that goal-directed cognition requires only 5\% additional energy above the brain's 20-watt baseline. This finding, typically cited to demonstrate the brain's efficiency, actually reveals something far more significant: \textbf{the 95\% baseline is not idle---it is active}.

That baseline metabolic expenditure maintains neural architecture, consolidates memory, and---crucially---provides capacity for the integrative processing that can generate genuine insight under appropriate conditions. \citet{wiehler2022} document the mechanism: sustained focused attention leads to glutamate accumulation in the lateral prefrontal cortex, creating the subjective experience of mental fatigue. The brain \textit{forces} disengagement from directed cognition not from weakness but from architectural necessity.

The implication restructures our understanding of learning. Synthesis emerges during apparent idleness---but not during \textit{any} form of idleness. Default mode network activation correlates with integrative processing, but the relationship is conditional. Low-threat-load states---contemplation without anxiety, reflection without deadline pressure---provide conditions for cross-domain pattern integration and recombination. High-threat-load ``rest''---rumination about economic precarity, vigilance during stressful commutes, passive consumption of attention-optimized media---activates similar neural regions but produces perseveration, stress consolidation, and attentional fragmentation rather than insight.

Historical educational systems, whether by design or accident, honored this architecture. Greek philosophical schools featured extended periods of protected contemplation. Medieval monasteries structured days around prayer, labor, and \textit{otium}---productive leisure in low-threat environments where survival was guaranteed. Renaissance patronage freed scholars from economic anxiety, enabling the \textit{protected} cognitive conditions that generated breakthrough synthesis.

Modern educational systems systematically eliminate this processing opportunity. The student juggling coursework, employment, commuting, and digital distraction has minimal unstructured cognitive space---and what remains is typically consumed by anxiety about the next deadline rather than low-threat-load integration. Their 5\% active cognition captures little because their 95\% baseline processes fragmentation and survival stress rather than consolidation and recombination.

\subsection{A Taxonomy of Cognitive Work}

Before examining demographic patterns, we must distinguish categories that contemporary statistics conflate. We propose a three-tier taxonomy:

\begin{enumerate}
    \item \textbf{Synthesizers}: Individuals whose primary work involves generating net-new conceptual frameworks, high-transfer insights, or navigation of genuine uncertainty. Historical examples include natural philosophers, theological innovators, and polymathic scholars. Modern proxies include research faculty, doctoral candidates engaged in original work, and R\&D personnel working at conceptual frontiers.
    
    \item \textbf{Knowledge practitioners}: Individuals performing analysis, design, coordination, or complex problem-solving within established frameworks. This category involves significant cognitive engagement but operates primarily within known parameters. Examples include engineers applying established methods, physicians following diagnostic protocols, and analysts working within defined domains.
    
    \item \textbf{Vector executors}: Individuals performing procedural application of established knowledge with minimal transfer or uncertainty navigation. Work is characterized by reproducible outputs within stable contexts.
\end{enumerate}

This taxonomy matters because contemporary employment statistics---particularly ``knowledge-intensive services'' classifications---aggregate all three tiers. The 42\% figure commonly cited for European knowledge work represents employment in knowledge-intensive service sectors \citep{eurostat2024}, a broad category encompassing synthesis, practice, and execution roles indiscriminately. Genuine synthesizers likely constitute a small fraction of this aggregate.

\subsection{The Demographic Inversion: From Concentration to Dilution}

With this taxonomy established, the civilizational analysis becomes most striking. The question is not merely total energy available, but energy investment \textit{per synthesizer}.

Pre-industrial Europe directed its entire educational surplus toward approximately 0.1--0.2\% of its population engaged in scholarly synthesis. A population of roughly 75 million at the pre-Black Death peak \citep{urlanis1941} concentrated its cognitive investment on fewer than 100,000 active scholars and university students. Contemporary estimates suggest Oxford never exceeded 1,600 enrolled students at any period; by 1300, approximately 23 universities operated across Europe, with Bologna---the largest---enrolling perhaps 500--1,500 students \citep{rashdall1936}. Adding monks engaged in scriptorial and scholarly work rather than agricultural labor yields a generous upper bound of 150,000 active synthesizers.

The percentage is striking: approximately 0.13\% of the pre-industrial European population engaged in synthesis-level intellectual work.

Modern Europe presents a structural inversion. Eurostat data indicate that 42.6\% of the employed European population now works in knowledge-intensive services, with 44.9\% classified as highly skilled workers (ISCO-08 major groups 1--3: managers, professionals, and technicians) \citep{eurostat2024}. From a population of approximately 450 million, this yields roughly 189 million employed persons in knowledge-intensive services---a broad category that includes many non-synthesis roles.

Even restricting to closer proxies for synthesis work---R\&D personnel, higher education researchers, doctoral candidates---yields perhaps 3--5 million Europeans engaged in activities approximating historical synthesis. This represents roughly 0.7--1.1\% of the population: a 5--8$\times$ expansion from pre-industrial levels, far less dramatic than the 320$\times$ expansion in knowledge-intensive employment overall.

The expansion in knowledge-intensive employment appears as democratization. Its thermodynamic consequence is dilution of synthesis-supporting conditions across a vastly larger population, most of whom occupy practitioner or executor roles.

\subsection{Externalized Survival: The Patronage Infrastructure}

Historical sphere-builders operated under conditions invisible to contemporary analysis: \textbf{externalized survival costs}. This was not incidental but structural.

\citet{holzknecht1923} documents the transformation: ``Medieval monasticism introduced a new element into the patronage of learning. It established a stable and sustainable model of institutional endowment that supported a lifetime of learning for generation after generation of monastics.'' Rather than individual scholars competing for patron attention---the precarious ancient model that left Seneca dependent on Nero's whims---medieval endowments provided unconditional maintenance. A lord ``endowed a monastery in perpetuity with a substantial gift of property,'' which ``afforded those with an interest in learning a secure position in a relatively well-endowed institution in which to pursue their studies comfortably.''

The thermodynamic significance is profound. Human cognitive capacity is finite. Energy devoted to survival---securing food, maintaining shelter, navigating economic systems, managing career anxiety---is energy unavailable for synthesis. Historical sphere-development required energy subsidies that freed cognitive resources for non-survival investment.

The trajectory can be characterized schematically, acknowledging that actual conditions varied widely across contexts and individuals. Pre-industrial scholars with fully externalized survival through endowments could allocate the substantial majority of their cognitive budget to scholarly activity. Pre-industrial artisans with integrated work and learning retained significant capacity for synthesis within their craft domains, though constrained by material production demands. Industrial workers with separated work and learning experienced reduction in available synthesis time as wage labor extracted focused hours. Modern students balancing employment, study, commuting, and debt service retain minimal cognitive surplus. Contemporary gig economy workers, fragmented across platforms with continuous economic uncertainty, approach conditions where synthesis becomes thermodynamically improbable.

The pattern is consistent: as survival costs are re-internalized onto individuals, cognitive surplus available for sphere-building collapses.

\subsection{The Per-Capita Calculation}

We can now quantify the inversion. The following calculations represent order-of-magnitude estimates intended to reveal structural patterns rather than precise measurements.

\textbf{Pre-industrial Europe (circa 1275):}

\begin{itemize}
    \item Population: approximately 75 million
    \item Active synthesizers: approximately 100,000 (0.13\%)
    \item Civilizational energy capacity: 75M $\times$ 20 GJ = 1.5 billion GJ/year
    \item Civilizational support capacity per synthesizer: 1.5B $\div$ 100k = \textbf{15,000 GJ/year}
\end{itemize}

This figure does not represent direct energy consumption by scholars. It represents civilizational investment capacity---the institutional overhead (buildings, libraries, food, heating, scribal materials, support staff) available per candidate synthesizer. The actual allocation to any individual varied enormously, but the \textit{capacity} for concentrated investment existed.

\textbf{Modern Europe (2024):}

\begin{itemize}
    \item Population: approximately 450 million
    \item Knowledge-intensive services employment: approximately 189 million (42\%)
    \item Civilizational energy capacity: 450M $\times$ 170 GJ = 76.5 billion GJ/year
    \item Energy capacity per KIS employee: 76.5B $\div$ 189M = \textbf{405 GJ/year}
\end{itemize}

Even using the more restrictive estimate of 4 million synthesis-proximate workers (R\&D, researchers, doctoral candidates), the figure becomes 76.5B $\div$ 4M = \textbf{19,125 GJ/year}---comparable to pre-industrial levels. However, this apparent parity obscures two critical differences: survival externalization and protected time.

\subsection{The Synthesis Investment Index}

We propose a heuristic index that makes explicit the factors determining synthesis-supporting conditions:

\begin{equation}
\text{SII} = \frac{\text{Civilizational energy capacity}}{\text{Candidate synthesizers}} \times \text{SE} \times \text{PT}
\label{eq:sii}
\end{equation}

Where:
\begin{itemize}
    \item \textbf{SII} = Synthesis Investment Index (heuristic units; not metabolic energy)
    \item \textbf{SE} (Survival Externalization) = fraction of cognitive capacity freed from earning subsistence (1.0 = fully endowed; 0.0 = fully self-funded)
    \item \textbf{PT} (Protected Time) = fraction of waking hours available for low-threat-load synthesis conditions
\end{itemize}

This decomposition is heuristic rather than physically rigorous---SE and PT are partially correlated with resource availability, and actual synthesis depends on factors beyond these three terms. The index captures structural conditions, not guaranteed outcomes.

\textbf{Pre-industrial scholars:}

Monastic horaria allocated 4--6 hours daily to study and copying, representing 40--60\% of productive waking hours \citep{knowles1963}. Whether individual monks achieved this in practice varied, but the institutional structure protected synthesis time in ways modern workplaces do not. Survival was unconditionally externalized through endowments.

\begin{equation}
\text{Pre-industrial SII: } 15,000 \times 1.0 \times 0.5 = \textbf{7,500 index units}
\end{equation}

\textbf{Modern knowledge-intensive workers:}

Contemporary research documents severe fragmentation. \citet{mark2008} found that office workers are interrupted or switch tasks every three minutes on average, with interruptions requiring an average of 23 minutes to resume original tasks. \citet{gonzalez2004} documented that information workers spend an average of only 11 minutes on any given project before switching. Industry surveys corroborate these patterns: knowledge workers report an average of 31.6 interruptions per day---approximately one every 15 minutes during focused work periods \citep{reclaimai2024}. Since achieving productive flow state requires 15--20 minutes of uninterrupted focus, this creates structural incompatibility with synthesis.

Cognitive peak capacity is limited to 2--4 hours per day \citep{ericsson1993}, yet knowledge workers spend substantial time in meetings, administrative tasks, and coordination activities that consume this limited capacity for non-synthesis purposes.

Survival is substantially internalized: knowledge workers must earn their maintenance through cognitive labor, redirecting capacity to economic survival even before accounting for administrative overhead. We estimate 10--20\% residual externalization through employer-provided benefits and infrastructure.

\begin{equation}
\text{Modern SII (KIS workers): } 405 \times 0.15 \times 0.05 = \textbf{3.0 index units}
\end{equation}

Even for synthesis-proximate workers with better structural conditions:
\begin{equation}
\text{Modern SII (researchers): } 19,125 \times 0.4 \times 0.15 = \textbf{1,147 index units}
\end{equation}

The ratio between pre-industrial and modern KIS workers is approximately 2,500$\times$. Between pre-industrial scholars and modern researchers, the ratio is approximately 6.5$\times$. These figures require careful interpretation: the calculation involves estimated parameters, and reasonable alternative assumptions yield different specific ratios. The structural finding---that modern conditions are substantially less favorable for synthesis than pre-industrial conditions, even controlling for population expansion---holds across reasonable parameter ranges.

\subsection{The Patronage Peak: A Conceptual Model}

Figure~\ref{fig:patronage-peak} presents a conceptual model of the relationship between knowledge worker density and peak synthesis capacity across historical periods.

\begin{figure}[htbp]
    \centering
    \includegraphics[width=0.9\textwidth]{patronage_peak.png}
    \caption{\textbf{The Patronage Peak (Conceptual Schematic).} This figure illustrates 
    the hypothesized relationship between knowledge worker percentage and peak 
    synthesis capacity across historical periods. The model proposes that maximum 
    synthesis capacity occurred when fewer than 2\% of the population engaged in 
    knowledge work but received concentrated patronage investment. As employment 
    in knowledge-intensive services expanded, per-capita investment fell below 
    sphere-building thresholds. \textit{Note: This is a conceptual illustration, not 
    an empirically derived curve. The vertical axis represents a theoretical construct 
    (synthesis capacity) that would require operationalization through proxies such as 
    polymathic output per capita or cross-domain innovation metrics. Specific values 
    are illustrative.}}
    \label{fig:patronage-peak}
\end{figure}

This pattern, if validated empirically, would demand careful interpretation. The claim is \textit{not} that education should be restricted to elites. The claim is thermodynamic: sphere-building requires sustained cognitive investment exceeding specific thresholds. When fixed educational resources distribute across expanding populations, per-capita investment may fall below those thresholds. The result would be universal access to sub-threshold education---democratized inadequacy.

\citet{turchin2023} identifies ``elite overproduction'' as a driver of civilizational instability: more credentialed candidates than available elite positions generates frustrated counter-elites who destabilize existing orders. Our framework suggests a deeper pathology. The problem may not be that we have overproduced elites; it is that we have eliminated the conditions for producing synthesizers. Post-Bologna reforms democratized credentials but may have eliminated the concentrated, decades-long investment that produced sphere-capable cognition. The solution would not be restricting access but \textbf{vertical extension}: democratize yesterday's elite tier while creating tomorrow's.

\subsection{The Energy Paradox: A Conceptual Resolution}

Figure~\ref{fig:energy-paradox} presents what appears paradoxical: humanity commands 10--50 times the per-capita energy of pre-industrial societies yet may produce substantially less cognitive synthesis per capita.

\begin{figure}[htbp]
    \centering
    \includegraphics[width=0.9\textwidth]{energy_paradox.png}
    \caption{\textbf{The Energy Paradox (Conceptual Schematic).} Civilizational energy 
    per capita (relative to pre-industrial baseline) plotted against estimated cognitive 
    surplus available for synthesis. Despite 10--50$\times$ energy increase since 1200 CE 
    \citep{smil2017}, the model proposes that protected synthesis time collapsed from 
    an estimated 40--60\% to approximately 3--5\%. The trajectory illustrates the 
    hypothesis that energy capture scaled substantially while extraction mechanisms 
    scaled faster still. \textit{Note: This is a conceptual illustration. Synthesis time 
    estimates for pre-industrial periods are inferred from institutional schedules 
    (e.g., monastic horaria) rather than direct measurement. Modern estimates derive 
    from time-use and interruption studies. Shaded region indicates uncertainty bounds.}}
    \label{fig:energy-paradox}
\end{figure}

The proposed resolution lies in \textbf{extraction scaling}. Energy capture increased substantially; mechanisms for extracting cognitive energy into non-synthesis channels may have scaled faster still.

Employment extraction captures 40--60 hours weekly of directed cognitive activity oriented toward organizational production rather than personal synthesis. Commute extraction consumes 5--15 hours weekly in transit that generates fatigue without integration opportunity. Administrative extraction demands continuous documentation, reporting, and compliance activities. Attention extraction, as \citet{zuboff2019} documents, deploys platform-mediated cognitive capture optimized through behavioral prediction. Anxiety extraction consumes background processing capacity through economic precarity and continuous career uncertainty.

Each extraction mechanism appeared locally rational---efficient employment, urban concentration, organizational accountability, engaging media, flexible labor markets. Collectively, they may constitute a civilizational system that approaches extraction saturation, leaving minimal residual for the synthesis that maintains adaptive capacity.

The thermodynamic implication is sobering: we may have optimized extraction while eliminating synthesis conditions. Pre-industrial societies with primitive energy technology produced Newton and Leibniz. Modern societies with abundant energy produce credentials. Whether this represents a fundamental trade-off or a correctable allocation failure remains an open question.

\subsection{The Modern Inversion}

The hidden energy budget framework suggests why contemporary attempts at sphere reconstruction systematically struggle. The conditions enabling historical sphere-building have been substantially inverted:

\begin{table}[htbp]
\centering
\caption{The modern inversion of sphere-building conditions}
\begin{tabular}{lll}
\toprule
\textbf{Factor} & \textbf{Historical} & \textbf{Contemporary} \\
\midrule
Survival costs & Externalized (endowments) & Largely internalized (wages) \\
Cognitive idleness & Structured, protected & Minimized, often stigmatized \\
Patronage/funding & Unconditional, lifetime & Conditional, project-based \\
Time horizon & Decades & Quarters to years \\
Economic anxiety & Absorbed by institutions & Significant background load \\
Attention environment & Low-stimulus & Optimized for capture \\
Task switching & Hours between contexts & Minutes between contexts \\
Social validation & Wisdom-based & Credential-based \\
\bottomrule
\end{tabular}
\label{tab:modern-inversion}
\end{table}

Multiple factors supporting sustained cognitive investment have been reversed. The modern knowledge worker operates in conditions potentially hostile to sphere development: survival concerns consuming baseline processing, employment extracting focused cognition, attention economy capturing residual awareness, credential systems rewarding vector optimization over synthesis capability.

This is not conspiracy but optimization. Each inversion emerged from locally rational decisions: internalize costs for efficiency, minimize idle time for productivity, condition funding on deliverables, shorten horizons for accountability, maintain urgency for motivation, capture attention for engagement, credential for standardization. The collective result---conditions making sphere-building difficult---was never intended but may now be substantially complete.

\subsection{Implications for Reconstruction}

The hidden energy budget framework yields specific requirements for any reconstruction attempt. Sphere-building appears to demand conditions that modern systems have systematically reduced.

Protected synthesis time requires a substantial threshold---we estimate 30--40\% of cognitive budget based on historical patterns---for unstructured, low-threat-load processing. This is not leisure but the conditions under which consolidation and recombination can occur. Modern knowledge workers achieving approximately 3--5\% protected time may operate below synthesis thresholds regardless of intelligence or motivation.

Externalized survival load requires that economic anxiety be absorbed institutionally rather than managed individually. The pre-industrial endowment model achieved this through perpetual property grants; modern salary models substantially do not, as workers must continuously earn their maintenance through cognitive labor. The cognitive overhead of earning survival---rent concern, healthcare uncertainty, retirement planning, career management---may consume the baseline processing that would otherwise support synthesis.

Extended time horizons require decades rather than semesters for cognitive architecture development. Sphere-building is slow because neural architecture modification is slow. Quarterly deliverables and annual reviews optimize for vector production, not synthesis development.

Attention sovereignty requires protection from optimized capture mechanisms. The attention economy has developed extraction technologies of unprecedented sophistication. Sphere-building may require insulation from these systems, not merely willpower to resist them.

Unconditional investment requires funding without deliverable extraction. Project-based funding with milestone requirements produces vectors---the cognitive architecture that generates predictable outputs. Synthesis emerges from exploration without predetermined destinations.

These requirements represent constraints rather than preferences. Systems violating them may produce vectors regardless of curriculum content, pedagogical innovation, or institutional intention. Whether such conditions can be recreated at civilizational scale---or whether sphere-building necessarily remains confined to protected niches---determines whether human cognitive sovereignty survives the current transition. Section~\ref{sec:reconstruction} examines reconstruction possibilities within these constraints.

% ============================================
% END SECTION 5
% ============================================
% 
% FIGURE REQUIREMENTS:
% - patronage_peak.png: Add "CONCEPTUAL SCHEMATIC" label; remove precise axis numbers
% - energy_paradox.png: Add "CONCEPTUAL SCHEMATIC" label; add uncertainty shading
%
% NEW CITATIONS ADDED (v3.0):
%
% @article{mark2008,
%   author = {Mark, Gloria and Gudith, Daniela and Klocke, Ulrich},
%   title = {The Cost of Interrupted Work: More Speed and Stress},
%   journal = {Proceedings of CHI 2008},
%   pages = {107--110},
%   year = {2008},
%   publisher = {ACM}
% }
%
% @article{gonzalez2004,
%   author = {Gonz{\'a}lez, Victor M. and Mark, Gloria},
%   title = {Constant, Constant, Multi-tasking Craziness: Managing Multiple Working Spheres},
%   journal = {Proceedings of CHI 2004},
%   pages = {113--120},
%   year = {2004},
%   publisher = {ACM}
% }
%
% @article{ericsson1993,
%   author = {Ericsson, K. Anders and Krampe, Ralf Th. and Tesch-R{\"o}mer, Clemens},
%   title = {The Role of Deliberate Practice in the Acquisition of Expert Performance},
%   journal = {Psychological Review},
%   volume = {100},
%   number = {3},
%   pages = {363--406},
%   year = {1993}
% }
%
% CITATIONS RETAINED FROM v2.0:
% - smil2017, dezeeuw1978, urlanis1941, rashdall1936
% - holzknecht1923, eurostat2024, knowles1963
% - jamadar2025, wiehler2022, turchin2023, zuboff2019
% - reclaimai2024 (now secondary to Mark et al.)
%
% KEY CHANGES FROM v2.0:
% 1. Added 3-tier taxonomy (synthesizers/practitioners/executors)
% 2. Changed "Synthesis Capacity" → "Synthesis Investment Index (SII)"
% 3. Softened "performs integrative processing" → "provides capacity for"
% 4. Added Mark (2008), González & Mark (2004), Ericsson (1993)
% 5. Reframed figures as "Conceptual Schematic"
% 6. Changed "Medieval produced Newton" → "Pre-industrial produced Newton"
% 7. Clarified "189M knowledge workers" → "189M in KIS employment"
% 8. Added synthesis-proximate worker calculation (researchers/R&D)
% 9. Softened absolute claims throughout ("may," "appears to," "suggests")
% ============================================